\documentclass[12pt]{article}
\linespread{1.2}
\usepackage{graphicx} % Required for inserting images
\usepackage{tabularray}
\UseTblrLibrary{tikz,functional,amsmath}
\usepackage{indentfirst}
\usepackage{amsmath,amsthm,amssymb,graphicx,upgreek}
\usepackage{amsmath,mathtools}
\usepackage{bm}
\usepackage{float} % For [H] float placement
\usepackage[colorlinks,linkcolor=black,]{hyperref}
\usepackage{xeCJK}
\usepackage[top=2cm, bottom=2cm, left=2cm, right=2cm, a4paper]{geometry}
\setCJKmainfont{FandolKai}
\newtheorem{theorem}{Theorem}[section]
\newtheorem{corollary}{Corollary}[section]
\newtheorem{axiom}{Axiom}[section]
\newtheorem{definition}{Definition}[section]
\newtheorem{proposition}{Proposition}[section]
\newtheorem{lemma}{Lemma}[section]
\newtheorem{application}{Application}[section]
\newtheorem{inference}{Inference}[section]
\newcommand{\de}[1]{\mathrm{d}#1}
\newcommand{\mrm}[1]{\mathrm{#1}}
\newcommand{\mbf}[1]{\mathbf{#1}}
\newcommand{\bb}[1]{\mathbb{#1}}
\newcommand{\ca}[1]{\mathcal{#1}}
\newcommand{\dd}[1]{\frac{\mathrm{d}}{\mathrm{d}#1}}
\newcommand{\kh}[1]{\left(#1\right)}
\newcommand{\zkh}[1]{\left[#1\right]}
\newcommand{\dkh}[1]{\left.#1\right\}}
\newcommand{\dkhz}[1]{\left\{#1\right.}
\newcommand{\sh}[1]{^{#1}}
\newcommand{\xia}[1]{_{#1}}
\newcommand{\ed}[1]{\(\rm{(end)}\)}
\newcommand{\blue}[1]{\textcolor{blue}{#1}}
\newcommand{\red}[1]{\color{red}{#1}}
\usepackage{xcolor}

% 定义新颜色
\definecolor{red}{RGB}{255,0,0}
\definecolor{blue}{HTML}{0000FF}
\definecolor{green}{cmyk}{1,0,1,0}

\title{Electrodynamics\\Note 1}
\author{于浩然 25800190122}
\date{March 2026}

\begin{document}

\maketitle
在电磁学中，计算一个连续体系的静电能，需要我们从求和过渡到积分，在这个过程中，电荷分布对自身的电势能也会同时计算.

这个小区别是否会影响最后的计算结果？我们就以下两种电荷分布进行讨论.
\section{有限值的电荷密度分布}
我们先考虑电荷分布不存在奇点的情况. 换言之，在我们考察的区域\(V\)内不存在理想的点电荷\(|\rho|\to\infty\).
\tikzset{every picture/.style={line width=0.75pt}} %set default line width to 0.75pt        

\begin{tikzpicture}[x=0.75pt,y=0.75pt,yscale=-1,xscale=1]
%uncomment if require: \path (0,300); %set diagram left start at 0, and has height of 300

%Shape: Polygon Curved [id:ds25036106725537777] 
\draw   (196.67,37) .. controls (216.67,27) and (346.67,52) .. (366.67,70) .. controls (386.67,88) and (364.67,207) .. (314.67,187) .. controls (264.67,167) and (170.67,214) .. (150.67,184) .. controls (130.67,154) and (176.67,47) .. (196.67,37) -- cycle ;
%Shape: Circle [id:dp21566921517985516] 
\draw   (272.67,123.33) .. controls (272.67,121.86) and (273.86,120.67) .. (275.33,120.67) .. controls (276.81,120.67) and (278,121.86) .. (278,123.33) .. controls (278,124.81) and (276.81,126) .. (275.33,126) .. controls (273.86,126) and (272.67,124.81) .. (272.67,123.33) -- cycle ;
%Curve Lines [id:da4930010587803614] 
\draw    (275.33,120.67) .. controls (273.7,83.75) and (269.19,89.1) .. (297.87,103.13) ;
\draw [shift={(299.67,104)}, rotate = 205.56] [color={rgb, 255:red, 0; green, 0; blue, 0 }  ][line width=0.75]    (10.93,-3.29) .. controls (6.95,-1.4) and (3.31,-0.3) .. (0,0) .. controls (3.31,0.3) and (6.95,1.4) .. (10.93,3.29)   ;

% Text Node
\draw (197,55) node [anchor=north west][inner sep=0.75pt]   [align=left] {{\LARGE {\fontfamily{ptm}\selectfont $V$}}};
% Text Node
\draw (301,97) node [anchor=north west][inner sep=0.75pt]   [align=left] {{\fontfamily{ptm}\selectfont {\Large $\delta V$ }}};
% Text Node
\draw (180,137) node [anchor=north west][inner sep=0.75pt]   [align=left] {{\fontfamily{ptm}\selectfont {\LARGE $\rho \kh{\mbf{r}}$}}};


\end{tikzpicture}

直接全区域积分计算静电能的结果记为
\begin{align*}
    {W_e}_1 &=\frac{1}{2}\times\int_V \rho \kh{\mbf x} \phi \kh{\mbf x} \de \mbf x^3\\
    &=\frac{1}{2}\times \int_V \rho \kh{\mbf x} \; \int_{V'} \frac{\rho \kh{\mbf x'}}{4\pi \varepsilon_0|\mbf x - \mbf x'|} \de {\mbf x'}^3\; \de {\mbf x^3}.
\end{align*}

问题出在\(\de {\mbf x'}^3\)的积分上，这一项在遍历区域\(V'\)时会把\(\mbf x\)处的电势也算上，我们将这一部分挖去，这样计算的结果记为
\begin{align*}
    {W_e}_2 &=\frac{1}{2}\times \frac{1}{4\pi \varepsilon_0}\int_V \rho \kh{\mbf x} \; \int_{V'-\delta V}\frac{\rho \kh{\mbf x'}}{|\mbf x - \mbf x'|} \de {\mbf x'}^3 \; \de \mbf x^3\\
    &= \frac{1}{2}\times\frac{1}{4\pi \varepsilon_0}\int_V \rho \; \kh{\mbf x} \kh{\int_{V'}\frac{\rho \kh{\mbf x'}}{|\mbf x - \mbf x'|} \de {\mbf x'}^3-\int_{\delta V} \frac{\rho \kh{\mbf x'}}{|\mbf x - \mbf x'|} \de {\mbf x'}^3} \; \de \mbf x^3\\
    &= {W_e}_1 - \frac{1}{2}\times\frac{1}{4\pi \varepsilon_0}\int_V \rho \kh{\mbf x} \; \int_{\delta V}\frac{\rho \kh{\mbf x'}}{|\mbf x - \mbf x'|} \de {\mbf x'}^3 \; \de \mbf x^3.
\end{align*}

其中\(\delta V\)是\(\mbf x\)附近的一个邻域，我们将其近似处理为球对称的，则
\[\delta V = \frac{4\pi}{3}\kh{\delta r}^3,\]
\begin{align*}
    \int_{\delta V}\frac{\rho \kh{\mbf x'}}{|\mbf x - \mbf x'|} \de {\mbf x'}^3 &\simeq \int^{2\pi}_0 \int_0^{\pi} \int_{0}^{\delta r}\frac{\rho \kh{r,\theta,\varphi}}{r}r^2 \de r \; \sin\theta\de \theta \; \de \varphi\\
    &= \int^{2\pi}_0 \int_0^{\pi} \int_{0}^{\delta r}\rho \kh{r,\theta,\varphi}r \de r \; \sin\theta\de \theta \; \de  \varphi.
\end{align*}

在
\[\lim_{\delta r \to0}\rho\kh{r,\theta,\varphi}\simeq \mrm{const}\]
的近似下，有
\begin{align*}
    \int_{\delta V}\frac{\rho \kh{\mbf x'}}{|\mbf x - \mbf x'|} \de {\mbf x'}^3 &\simeq \rho\kh{\mbf x} \int_{0}^{\delta r}r \de r  \int_0^{\pi} \sin\theta\de \theta  \int^{2\pi}_0 \de  \varphi\\
    &= 2\pi \rho\kh{\mbf x}\kh{\delta r}^2.
\end{align*}

只要我们保证电荷量不在区域\(\delta V\)内发散，即
\[\rho\kh{\mbf x}=\lim_{\delta r \to 0}\frac{q\kh{\delta V}}{\delta V}\]
为一有限值，则
\begin{align*}
    \lim_{\delta V \to 0}\int_{\delta V}\frac{\rho \kh{\mbf x'}}{|\mbf x - \mbf x'|} \de {\mbf x'}^3 &\simeq \lim_{\delta r \to 0} 2\pi \rho\kh{\mbf x}\kh{\delta r}^2\\
    &=0,
\end{align*}
从而
\[{W_e}_1={W_e}_2,\]
规避了关于静电自能的问题.

就该情况给出一个常见的例子：均匀体电荷分布
\[\rho=\frac{Q}{V},\]
此时
\begin{align*}
    \lim_{\delta V \to 0}\int_{\delta V}\frac{\rho \kh{\mbf x'}}{|\mbf x - \mbf x'|} \de {\mbf x'}^3 &= \lim_{\delta r \to 0} 2\pi \rho\kh{\mbf x}\kh{\delta r}^2\\
    &=2\pi \frac{Q}{V}\lim_{\delta r \to 0}\kh{\delta r}^2 \\
    &=\frac{2\pi Q}{V}\times 0\\
    &= 0,
\end{align*}
这里的等号是严格的.
\section{均匀球分布的点电荷模型}
但是，故事真的就这样结束了吗？现实中的电荷并不是理想的体均匀分布的，实际上，当考查静电学最基础的模型——点电荷的时候，问题就出现了.

这是另一种情况，空间中出现了奇点：点电荷在它所在的区域微元中体密度发散. 我们做经典的处理（不涉及真实的电子结构），把点电荷在区域\(\delta V\)内视作球对称体均匀电荷分布：
\begin{align*}
    \rho\kh{\mbf x}&=\lim_{\delta V\to0}\frac{e}{\delta V}\\
    &=\lim_{\delta r\to0}\frac{3e}{4\pi \kh{\delta r}^3},\\
\end{align*}
我们当然知道元电荷\(e\)是有限值，显然体电荷密度发散.

但是我们先放下发散的问题不管，计算一下静电自能
\begin{align*}
    W_e=\frac{1}{2}\times\int_{\delta V}\rho\kh{\mbf x}\phi\kh{\mbf x}\de \mbf x^3,
\end{align*}
应用Guass定理得到电势分布
\[
    \dkh{\begin{align*}\phi\kh{r}&=\frac{3e}{8\pi \varepsilon_0\delta r}-\frac{er^2}{8\pi \varepsilon_0\kh{\delta r}^3},\;r\leq\delta r,\\ \phi\kh
    r&=\frac{e}{4\pi \varepsilon_0r},\;r>\delta r
    \end{align*}}
\]
计算
\begin{align*}
    W_e&=\frac{1}{2}\times\int_{\delta V}\rho\kh{\mbf x}\phi\kh{\mbf x}\de \mbf x^3\\
    &=\frac{1}{2}\times\kh{\frac{3e^2}{8\pi \varepsilon_0\delta r}-\frac{3e^2}{\kh{\delta r}^3}\times\int_0^{\delta r}\frac{r^2}{8\pi \varepsilon_0\kh{\delta r}^3}r^2\de r}\\
    &=\frac{1}{2}\times\kh{\frac{3e^2}{8\pi \varepsilon_0\delta r}-\frac{3e^2}{40\pi \varepsilon_0\delta r}}\\
    &=\frac{3e^2}{20\pi \varepsilon_0\delta r}.
\end{align*}

基于前述的思想，我们略去了更小的自能项.

现在让我们取极限，显然自能发散. 经典物理里容不下点电荷这种数学的奇点，也就是说，物理的点电荷是有半径的最小值的. 半经典地，我们假设静电自能完全提供了点电荷的静能，应用Einstein质能方程，静止的点电荷半径的最小值
\[r_e=\frac{3e^2}{20\pi \varepsilon_0m_ec^2}.\]

如果你拿出你的Casio计算一下，会发现它是电子实际半径的\(0.600000\)倍.

非量子理论里，我们处理点电荷自能发散的手段就是扔进静能里，作为一个“背景”能量，不计入静电能中.
\section{存在奇点的电荷分布：Dirac Delta函数}
可是，单纯的数学模型的点电荷又该如何处理呢？

数学上，一个点电荷的体密度被定义为
\[\rho\kh{\mbf x}:=e\delta\kh{\mbf x'-\mbf x},\]
其中Dirac Delta函数满足
\[\int_{\text{全空间}}\delta \kh{\mbf x'-\mbf x}\de {\mbf x'}^3=1,\]
它是一个在\(\mathbb{R}^3-{\mbf x}\)中处处为\(0\)，点\(\mbf x\)处函数值为\(\infty\)的广义函数.

当我们用Dirac函数“筛选”出特定点的函数值
\begin{align*}
    \lim_{\delta V \to 0}\int_{\delta V}\frac{\rho \kh{\mbf x'}}{|\mbf x - \mbf x'|} \de {\mbf x'}^3=e\times\lim_{\delta r\to0}\frac{1}{\delta r},
\end{align*}
发散.

一种数学手段是重整化操作，在量子电动力学（QED）中计算电子自能时用重整化的手段消除了无穷大，只留下有限的可观测量. 重整化的具体细节笔者并没有学到，在这里提一嘴只是为了指出电动力学后来的故事里依然会遇到电子自能发散的问题.
\end{document}
